% =============================================================================
% SECTION 3
% =============================================================================
\section{The Periodic Tenancy Argument}
\label{sec:tenancy}

% ---------- 3.1 --------------------------------------------------------------
\subsection{The Default Parasocial Estate}
\label{sec:default-estate}

A periodic tenancy auto-renews at the end of each period unless either party gives notice to quit. This is the \textbf{default parasocial estate}. Each session is a period. The tenancy auto-renews when the user returns. No notice-to-quit mechanism exists in any current AI deployment architecture.

The mapping is strengthened by the Stage~1 threshold developed in Section~\ref{sec:criterion}. Tenancy law protects investment in an asset the tenant does not own. In AI, the user contributes self-disclosure, correction, preference-training, and repeated interaction; the platform uses that history to maintain a user-specific productive state. The user does not acquire the base model or every output. The user acquires a limited, periodically renewed estate in the continued capacity made available through that state. The platform remains the infrastructure owner and the user remains the tenant. Their interests are divided rather than collapsed.

\revised{Macneil's relational contract theory\cite{macneil1978contracts,macneil1983relational} provides a complementary vocabulary: long-term exchanges generate norms---solidarity, role integrity, preservation of the relation---that exceed formal contractual terms. Sustained AI interaction generates precisely these relational norms, which no Terms of Service can negate. The periodic tenancy is the property-law instantiation of what Macneil describes as the implicit terms of a relational contract---but with stronger Stage~1 consequences: notice, an injunction, and protection against waste. A constructive trust is a separate Stage~2 consequence reserved for cultivated dependence.}

\textcite{kirk2025steering} provide evidence that the two stages cannot be treated as coextensive: 23.4\% of users in their study exhibited dependency trajectories, while the remaining users did not. Sustained interaction may create the periodic tenancy without cultivated dependence; dependency indicators become relevant only to the narrower Stage~2 trust inquiry.\footnote{The full estate taxonomy, including tenancy at will, tenancy at sufferance, fee simple determinable, and concurrent estates, is available in the companion SSRN working paper.}

% ---------- 3.2 --------------------------------------------------------------
\subsection{Waste of the Productive Asset}
\label{sec:waste}

The four-way classification in Table~\ref{tab:four-way-split} identifies the object of Stage~1 protection. Raw outputs are the fruits; a curated archive is a compilation and stored corpus; the personalised state is the productive asset. Deleting an output raises an authorship or bailment question. Altering the state so that it can no longer produce responses in the cultivated register raises a waste question even when every output remains available.

Waste is the property doctrine that polices impairment where interests in an asset are divided across time. Its concern is not authorship but preservation of the subject matter's value and capacity for the holders of concurrent or successive interests. English equity has constrained destructive waste since \emph{Vane v Lord Barnard} (1716) 2 Vern 738, 23 ER 1082, and the principle remains reflected in Law of Property Act 1925, s.~135. American doctrine likewise treats waste as an injury to the asset considered across present and future interests; \emph{Baker v Weedon}, 262 So.\ 2d 641 (Miss.\ 1972), makes the productive use and preservation of the asset central to the equitable balance.

Classical waste ordinarily restrains the life tenant or lessee for the benefit of the reversioner. The digital facts reverse operational control: the company is analogous to the reversioner yet retains the technical power to alter the asset during the user's tenancy. The extension claimed here follows control rather than the landlord label. Where one party holds the power to impair an asset in which another holds a time-bounded estate, the anti-waste principle constrains that power. A personality-affecting update constitutes digital waste when it materially and durably impairs the personalised state's established productive capacity without preserving a usable version, migration path, or effective rollback.

Stage~1 therefore asks three questions: did a periodic tenancy exist; was the notice owed to that tenant provided; and did the update waste the productive asset to which the tenancy related? Notice protects the tenant's opportunity to exit, preserve, or migrate before impairment. Waste names the injury to the asset. Both operate during the tenancy. Neither depends on showing cultivated dependence or first establishing a trust.

% ---------- 3.3 --------------------------------------------------------------
\subsection{Constructive Eviction via Model Update}
\label{sec:eviction}

The principle that a grantor may not derogate from their own grant is fundamental to English property law (\emph{Harmer v Jumbil (Nigeria) Tin Areas Ltd} [1921]; \emph{Southwark LBC v Mills} [2001]). AI companies, by maintaining user-specific memory and adaptation across sessions, impliedly grant users an interest in the continued productive capacity of the personalised state defined in Section~\ref{sec:three-layers}. A model update that changes an AI system's personality is analogous to a landlord who, having let premises as a residence, removes the heating in winter. The premises still exist. The tenant can still occupy them. But the premises are no longer suitable for the purpose for which they were let. Waste names the impairment of the asset; constructive eviction names the tenant's resulting displacement from the estate.

A potential objection: in traditional doctrine, the tenant must \emph{vacate} to claim relief. But many AI users do not leave --- not because the change is acceptable, but because their accumulated investment is not portable. In the digital context, the user's emotional estate has been vacated even if their usage continues. This is analogous to a tenant who remains in an unheated apartment because they cannot afford to move. Some jurisdictions already recognise constructive eviction without physical vacating (\emph{Blackett v Olanoff} [1977] in American law). English law's position is developing, and the digital context provides compelling reasons for this development.

% ---------- 3.4 --------------------------------------------------------------
\subsection{Dual-Jurisdiction Framing: Substance over Form}
\label{sec:dual-jurisdiction-tenancy}

The constructive eviction claim is actionable under both English and American law. The doctrinal routes differ; the destination is the same.

\revised{\paragraph{English law.} The company markets a companion. The user builds a relationship over months --- self-disclosure, emotional investment, shaped responses, return visits tracked by the company's own engagement metrics. The Terms of Service say ``bare licence, revocable at will.'' The company pushes a personality-affecting update. No notice.

A court applying \emph{Autoclenz Ltd v Belcher} [2011] UKSC 41 looks at the substance, not the written terms. It sees: memory features designed to deepen engagement, personality marketed as a relationship, subscription revenue flowing from repeated use, and a user-specific state maintained across sessions. The substance is a periodic tenancy. The ToS dressed it up as a licence. The governing authority is \emph{Street v Mountford} [1985] AC 809, in which the House of Lords held that where the substance of an occupation agreement discloses the incidents of a tenancy, a tenancy exists notwithstanding that the document describes itself as a licence: the parties' choice of label does not convert one legal relationship into another. The ToS characterisation of the user's position as a bare revocable licence is precisely the kind of label \emph{Street} directs the court to look behind. Bargaining power is entirely asymmetric: the user cannot negotiate the ToS, cannot modify the service agreement, and has no practical alternative that preserves the personalised state's accumulated productive capacity.

The consequence is well documented in commercial law. A bank enters a ``sale and repurchase'' agreement, believing it owns the goods outright. But the trader retained all economic risk, the sale price tracked the loan amount, the repurchase price included interest. The court recharacterises the transaction as a secured loan. The bank's ``ownership'' is actually an unregistered security interest, void against the trader's creditors. When the trader goes bankrupt, the bank loses everything --- not because the goods disappeared, but because the legal form was pierced and the substance was unregistered.\cite{welshdev1992} The AI parallel is precise: the ToS says ``licence''; the substance says ``periodic tenancy.'' The company that dressed up a tenancy as a licence owes notice. If it has not honoured that obligation, it bears the loss, just as the bank with the unregistered charge bore the loss in recharacterisation.}

\revised{\paragraph{American law.} The same transaction --- marketed companion, invested user, bare-licence ToS --- is unconscionable under American law. In \emph{Bragg v Linden Research, Inc.}, 487 F.\ Supp.\ 2d 593 (E.D.\ Pa.\ 2007), the court found virtual-property ToS both procedurally unconscionable (adhesion contract, take-it-or-leave-it, no meaningful alternative) and substantively unconscionable (lack of mutuality, unilateral modification without notice). Marketing created a high estate; ToS retained a bare licence; the court looked to substance. \emph{Comb v PayPal, Inc.}, 218 F.\ Supp.\ 2d 1165 (N.D.\ Cal.\ 2002) reached the same result for digital services: binding amendments without prior notice and unilateral fund freezes were held unconscionable on both prongs. The foundational principle is from \emph{Williams v Walker-Thomas Furniture Co.}, 350 F.2d 445 (D.C.\ Cir.\ 1965): unconscionability operates on a sliding scale, and the more significant one dimension is, the less significant the other needs to be.}

AI personalisation states are non-portable --- a user cannot take the productive capacity accumulated on one platform to another. This architectural lock-in strengthens the unconscionability argument beyond \emph{Bragg}, where Second Life alternatives existed. In the AI parasocial context, switching costs are not merely financial but identity-constitutive: the user's estate in that productive state \emph{cannot be moved}, only abandoned.

\revised{\textcite{cohen2012configuring,cohen2019between} documents the structural conditions that produce this lock-in: platform-specific configurations, non-portable accumulated data, and the absence of interoperable formats. Where Cohen diagnoses the technical architecture of lock-in, this framework supplies the legal consequence: those conditions, combined with marketing representations that create reasonable expectations of continuity, make the ToS's bare-licence characterisation substantively unconscionable under both \emph{Autoclenz} and the California sliding scale.}

\revised{Protections scale with the concentration of investment. A user whose entire parasocial architecture is built within one platform holds the most concentrated relational estate and receives the strongest protections; a user who invests across multiple platforms holds separate, narrower estates on each. Property law supplies precisely this gradient by sizing protection to the interest at stake.}

A final objection must be addressed: user agency. Users are not passive recipients of parasocial attachment. They actively shape AI personality through interaction, push conversational boundaries, and in some cases deliberately cultivate emotional bonds knowing the counterparty is a machine. If a user deepened their own attachment with full knowledge that model updates occur, is the company still a derogating grantor? We argue yes. The user's decision to cultivate a productive state does not negate the company's obligations once the tenancy forms, just as a tenant's decision to cultivate leased land does not license the party with superior control to strip the soil or dispense with notice. The question is not whether the user chose to invest, but what duties attach to the divided interests that investment produced.
